% Priority scheduling of user-level threads on two CPUs: when T2 unlocks the
% mutex, the library must signal the other kernel-level thread to preempt T1.
% Usage: % Priority scheduling of user-level threads on two CPUs: when T2 unlocks the
% mutex, the library must signal the other kernel-level thread to preempt T1.
% Usage: % Priority scheduling of user-level threads on two CPUs: when T2 unlocks the
% mutex, the library must signal the other kernel-level thread to preempt T1.
% Usage: % Priority scheduling of user-level threads on two CPUs: when T2 unlocks the
% mutex, the library must signal the other kernel-level thread to preempt T1.
% Usage: \input{figures/l05-multi-cpu}   (width ~116 mm)
\begin{tikzpicture}[x=1mm, y=1mm, font=\footnotesize\sffamily,
  lab/.style={font=\scriptsize\sffamily, align=center},
  row/.style={font=\scriptsize\sffamily, anchor=east, align=right},
  >={Stealth[length=1.6mm]}]
  % \lfseg{x1}{x2}{y}{fill}{text}
  \def\lfseg#1#2#3#4#5{%
    \draw[fill=#4] (#1,#3-2.5) rectangle (#2,#3+2.5);
    \node[lab] at ({(#1+#2)/2},#3) {#5};}
  % event lines
  \draw[densely dashed, ln-muted] (44,24) -- (44,-2);
  \draw[densely dashed, ln-muted] (74,24) -- (74,-2);
  \node[lab, anchor=south, fill=white, inner sep=1pt] at (44,24)
    {\iflangja{T3 が m でブロック}{T3 blocks on m}};
  \node[lab, anchor=south, fill=white, inner sep=1pt] at (74,24)
    {\iflangja{T2 が m を解放}{T2 unlocks m}};
  % CPU 1
  \node[row] at (24,15) {CPU 1\\(KLT K1)};
  \lfseg{26}{74}{15}{ln-box}{\iflangja{T2（m を保持）}{T2 (holds m)}}
  \lfseg{74}{116}{15}{white}{T2}
  % CPU 2
  \node[row] at (24,4) {CPU 2\\(KLT K2)};
  \lfseg{26}{44}{4}{white}{T3}
  \lfseg{44}{48}{4}{black!35}{}
  \lfseg{48}{79}{4}{white}{T1}
  \lfseg{79}{83}{4}{black!35}{}
  \lfseg{83}{116}{4}{ln-box}{\iflangja{T3（m を保持）}{T3 (holds m)}}
  % signal from K1 to K2
  \draw[->, thick, ln-caution] (74,12.5) -- node[lab, right, pos=0.45, text=ln-caution]
    {\iflangja{シグナル}{signal}} (79,6.5);
  % legend
  \draw[fill=black!35] (26,-9) rectangle (30,-6);
  \node[lab, anchor=west] at (31,-7.5) {\iflangja{ライブラリのスケジューラ}{library scheduler}};
  \draw[fill=ln-box] (64,-9) rectangle (68,-6);
  \node[lab, anchor=west] at (69,-7.5) {\iflangja{m を保持}{m held}};
  \draw[->] (92,-7.5) -- node[lab, above] {\iflangja{時間}{time}} (116,-7.5);
\end{tikzpicture}
   (width ~116 mm)
\begin{tikzpicture}[x=1mm, y=1mm, font=\footnotesize\sffamily,
  lab/.style={font=\scriptsize\sffamily, align=center},
  row/.style={font=\scriptsize\sffamily, anchor=east, align=right},
  >={Stealth[length=1.6mm]}]
  % \lfseg{x1}{x2}{y}{fill}{text}
  \def\lfseg#1#2#3#4#5{%
    \draw[fill=#4] (#1,#3-2.5) rectangle (#2,#3+2.5);
    \node[lab] at ({(#1+#2)/2},#3) {#5};}
  % event lines
  \draw[densely dashed, ln-muted] (44,24) -- (44,-2);
  \draw[densely dashed, ln-muted] (74,24) -- (74,-2);
  \node[lab, anchor=south, fill=white, inner sep=1pt] at (44,24)
    {\iflangja{T3 が m でブロック}{T3 blocks on m}};
  \node[lab, anchor=south, fill=white, inner sep=1pt] at (74,24)
    {\iflangja{T2 が m を解放}{T2 unlocks m}};
  % CPU 1
  \node[row] at (24,15) {CPU 1\\(KLT K1)};
  \lfseg{26}{74}{15}{ln-box}{\iflangja{T2（m を保持）}{T2 (holds m)}}
  \lfseg{74}{116}{15}{white}{T2}
  % CPU 2
  \node[row] at (24,4) {CPU 2\\(KLT K2)};
  \lfseg{26}{44}{4}{white}{T3}
  \lfseg{44}{48}{4}{black!35}{}
  \lfseg{48}{79}{4}{white}{T1}
  \lfseg{79}{83}{4}{black!35}{}
  \lfseg{83}{116}{4}{ln-box}{\iflangja{T3（m を保持）}{T3 (holds m)}}
  % signal from K1 to K2
  \draw[->, thick, ln-caution] (74,12.5) -- node[lab, right, pos=0.45, text=ln-caution]
    {\iflangja{シグナル}{signal}} (79,6.5);
  % legend
  \draw[fill=black!35] (26,-9) rectangle (30,-6);
  \node[lab, anchor=west] at (31,-7.5) {\iflangja{ライブラリのスケジューラ}{library scheduler}};
  \draw[fill=ln-box] (64,-9) rectangle (68,-6);
  \node[lab, anchor=west] at (69,-7.5) {\iflangja{m を保持}{m held}};
  \draw[->] (92,-7.5) -- node[lab, above] {\iflangja{時間}{time}} (116,-7.5);
\end{tikzpicture}
   (width ~116 mm)
\begin{tikzpicture}[x=1mm, y=1mm, font=\footnotesize\sffamily,
  lab/.style={font=\scriptsize\sffamily, align=center},
  row/.style={font=\scriptsize\sffamily, anchor=east, align=right},
  >={Stealth[length=1.6mm]}]
  % \lfseg{x1}{x2}{y}{fill}{text}
  \def\lfseg#1#2#3#4#5{%
    \draw[fill=#4] (#1,#3-2.5) rectangle (#2,#3+2.5);
    \node[lab] at ({(#1+#2)/2},#3) {#5};}
  % event lines
  \draw[densely dashed, ln-muted] (44,24) -- (44,-2);
  \draw[densely dashed, ln-muted] (74,24) -- (74,-2);
  \node[lab, anchor=south, fill=white, inner sep=1pt] at (44,24)
    {\iflangja{T3 が m でブロック}{T3 blocks on m}};
  \node[lab, anchor=south, fill=white, inner sep=1pt] at (74,24)
    {\iflangja{T2 が m を解放}{T2 unlocks m}};
  % CPU 1
  \node[row] at (24,15) {CPU 1\\(KLT K1)};
  \lfseg{26}{74}{15}{ln-box}{\iflangja{T2（m を保持）}{T2 (holds m)}}
  \lfseg{74}{116}{15}{white}{T2}
  % CPU 2
  \node[row] at (24,4) {CPU 2\\(KLT K2)};
  \lfseg{26}{44}{4}{white}{T3}
  \lfseg{44}{48}{4}{black!35}{}
  \lfseg{48}{79}{4}{white}{T1}
  \lfseg{79}{83}{4}{black!35}{}
  \lfseg{83}{116}{4}{ln-box}{\iflangja{T3（m を保持）}{T3 (holds m)}}
  % signal from K1 to K2
  \draw[->, thick, ln-caution] (74,12.5) -- node[lab, right, pos=0.45, text=ln-caution]
    {\iflangja{シグナル}{signal}} (79,6.5);
  % legend
  \draw[fill=black!35] (26,-9) rectangle (30,-6);
  \node[lab, anchor=west] at (31,-7.5) {\iflangja{ライブラリのスケジューラ}{library scheduler}};
  \draw[fill=ln-box] (64,-9) rectangle (68,-6);
  \node[lab, anchor=west] at (69,-7.5) {\iflangja{m を保持}{m held}};
  \draw[->] (92,-7.5) -- node[lab, above] {\iflangja{時間}{time}} (116,-7.5);
\end{tikzpicture}
   (width ~116 mm)
\begin{tikzpicture}[x=1mm, y=1mm, font=\footnotesize\sffamily,
  lab/.style={font=\scriptsize\sffamily, align=center},
  row/.style={font=\scriptsize\sffamily, anchor=east, align=right},
  >={Stealth[length=1.6mm]}]
  % \lfseg{x1}{x2}{y}{fill}{text}
  \def\lfseg#1#2#3#4#5{%
    \draw[fill=#4] (#1,#3-2.5) rectangle (#2,#3+2.5);
    \node[lab] at ({(#1+#2)/2},#3) {#5};}
  % event lines
  \draw[densely dashed, ln-muted] (44,24) -- (44,-2);
  \draw[densely dashed, ln-muted] (74,24) -- (74,-2);
  \node[lab, anchor=south, fill=white, inner sep=1pt] at (44,24)
    {\iflangja{T3 が m でブロック}{T3 blocks on m}};
  \node[lab, anchor=south, fill=white, inner sep=1pt] at (74,24)
    {\iflangja{T2 が m を解放}{T2 unlocks m}};
  % CPU 1
  \node[row] at (24,15) {CPU 1\\(KLT K1)};
  \lfseg{26}{74}{15}{ln-box}{\iflangja{T2（m を保持）}{T2 (holds m)}}
  \lfseg{74}{116}{15}{white}{T2}
  % CPU 2
  \node[row] at (24,4) {CPU 2\\(KLT K2)};
  \lfseg{26}{44}{4}{white}{T3}
  \lfseg{44}{48}{4}{black!35}{}
  \lfseg{48}{79}{4}{white}{T1}
  \lfseg{79}{83}{4}{black!35}{}
  \lfseg{83}{116}{4}{ln-box}{\iflangja{T3（m を保持）}{T3 (holds m)}}
  % signal from K1 to K2
  \draw[->, thick, ln-caution] (74,12.5) -- node[lab, right, pos=0.45, text=ln-caution]
    {\iflangja{シグナル}{signal}} (79,6.5);
  % legend
  \draw[fill=black!35] (26,-9) rectangle (30,-6);
  \node[lab, anchor=west] at (31,-7.5) {\iflangja{ライブラリのスケジューラ}{library scheduler}};
  \draw[fill=ln-box] (64,-9) rectangle (68,-6);
  \node[lab, anchor=west] at (69,-7.5) {\iflangja{m を保持}{m held}};
  \draw[->] (92,-7.5) -- node[lab, above] {\iflangja{時間}{time}} (116,-7.5);
\end{tikzpicture}
